\documentclass[11pt,a4paper]{article}

\usepackage[utf8]{inputenc}
\usepackage[T1]{fontenc}
\usepackage{lmodern}
\usepackage[margin=1in]{geometry}
\usepackage{amsmath,amssymb}
\usepackage{enumitem}
\usepackage{booktabs}
\usepackage{fancyhdr}
\usepackage{framed}
\usepackage{xcolor}

% Gray solution boxes
\definecolor{solngray}{gray}{0.95}
\newenvironment{solution}{\begin{shaded}\textbf{SOLUTION:}\par\smallskip}{\end{shaded}}
\newenvironment{shaded}{\begin{snugshade}}{\end{snugshade}}
\surroundwithmdframed[backgroundcolor=solngray,linewidth=0pt]{shaded}

% Actually use mdframed
\usepackage{mdframed}
\newmdenv[backgroundcolor=solngray,linewidth=0.5pt,linecolor=gray,innertopmargin=8pt,innerbottommargin=8pt]{solnbox}

\pagestyle{fancy}
\fancyhf{}
\fancyhead[L]{\textbf{Discrete Structures I}}
\fancyhead[R]{\textbf{Prelim Practice Exam A}}
\fancyfoot[C]{\thepage}
\renewcommand{\headrulewidth}{0.5pt}

\begin{document}

\begin{center}
{\LARGE\textbf{DISCRETE STRUCTURES I}}\\[5pt]
{\Large Prelim Practice Exam --- Version A}\\[5pt]
{\large\textit{With Complete Solutions}}\\[10pt]
\rule{0.8\textwidth}{0.5pt}\\[5pt]
\textbf{Time Allowed:} 1.5 hours \quad | \quad \textbf{Total Points:} 55
\end{center}

\vspace{10pt}

\noindent\textbf{Instructions:}
\begin{itemize}[noitemsep]
    \item Answer all questions.
    \item Show all work for full credit.
    \item Write clearly and box your final answers.
\end{itemize}

\vspace{15pt}

%==============================================================================
% PART I: MULTIPLE CHOICE
%==============================================================================
\section*{PART I: Multiple Choice (10 points)}
\textit{Circle the letter of the best answer. 2 points each.}

\vspace{8pt}

\textbf{1.} Which of the following is a proposition?

\begin{enumerate}[label=(\Alph*),noitemsep]
    \item $x + 5 = 10$
    \item What time is it?
    \item Close the door.
    \item $2 + 2 = 5$
\end{enumerate}

\begin{solnbox}
\textbf{Answer: (D)}

\textbf{Explanation:} A proposition is a declarative sentence that is either TRUE or FALSE, but not both.
\begin{itemize}[noitemsep]
    \item (A) $x + 5 = 10$ --- Not a proposition because it contains a variable $x$. We cannot determine its truth value without knowing $x$.
    \item (B) ``What time is it?'' --- This is a question, not a declarative statement.
    \item (C) ``Close the door.'' --- This is a command/imperative, not declarative.
    \item (D) $2 + 2 = 5$ --- This IS a proposition! It's a declarative statement that is definitively FALSE. Remember: a proposition doesn't have to be true, it just needs to have a definite truth value.
\end{itemize}
\end{solnbox}

\vspace{10pt}

\textbf{2.} The contrapositive of $p \to q$ is:

\begin{enumerate}[label=(\Alph*),noitemsep]
    \item $q \to p$
    \item $\neg p \to \neg q$
    \item $\neg q \to \neg p$
    \item $\neg p \lor q$
\end{enumerate}

\begin{solnbox}
\textbf{Answer: (C)}

\textbf{Explanation:} For any conditional $p \to q$:
\begin{itemize}[noitemsep]
    \item \textbf{Converse:} $q \to p$ (swap the order) --- NOT equivalent to original
    \item \textbf{Inverse:} $\neg p \to \neg q$ (negate both) --- NOT equivalent to original
    \item \textbf{Contrapositive:} $\neg q \to \neg p$ (swap AND negate both) --- IS equivalent to original!
\end{itemize}

\textbf{Memory trick:} The contrapositive is the ONLY related statement that is logically equivalent to the original. Think of it as ``if not the consequence, then not the cause.''
\end{solnbox}

\vspace{10pt}

\textbf{3.} What is the truth value of $\forall x (x^2 \geq 0)$ if the domain is all real numbers?

\begin{enumerate}[label=(\Alph*),noitemsep]
    \item True
    \item False
    \item Cannot be determined
    \item Depends on $x$
\end{enumerate}

\begin{solnbox}
\textbf{Answer: (A)}

\textbf{Explanation:} The statement says ``For ALL real numbers $x$, $x^2 \geq 0$.''

This is TRUE because:
\begin{itemize}[noitemsep]
    \item If $x > 0$: $x^2 > 0 \geq 0$ \checkmark
    \item If $x = 0$: $x^2 = 0 \geq 0$ \checkmark
    \item If $x < 0$: $x^2 > 0 \geq 0$ \checkmark (negative times negative = positive)
\end{itemize}

Since the square of any real number is never negative, this universal statement holds for the entire domain.
\end{solnbox}

\vspace{10pt}

\textbf{4.} Which rule of inference is represented by: ``If $p \to q$ and $\neg q$, then $\neg p$''?

\begin{enumerate}[label=(\Alph*),noitemsep]
    \item Modus Ponens
    \item Modus Tollens
    \item Hypothetical Syllogism
    \item Disjunctive Syllogism
\end{enumerate}

\begin{solnbox}
\textbf{Answer: (B)}

\textbf{Explanation:} 
\begin{itemize}[noitemsep]
    \item \textbf{Modus Ponens:} $p \to q$, $p$ $\therefore q$ (``If P then Q; P is true; so Q'')
    \item \textbf{Modus Tollens:} $p \to q$, $\neg q$ $\therefore \neg p$ (``If P then Q; Q is false; so P must be false'')
    \item \textbf{Hypothetical Syllogism:} $p \to q$, $q \to r$ $\therefore p \to r$ (chain implications)
    \item \textbf{Disjunctive Syllogism:} $p \lor q$, $\neg p$ $\therefore q$ (process of elimination)
\end{itemize}

The given form matches \textbf{Modus Tollens} --- we reason backwards from the consequence being false.
\end{solnbox}

\vspace{10pt}

\textbf{5.} The negation of $\exists x\, P(x)$ is:

\begin{enumerate}[label=(\Alph*),noitemsep]
    \item $\exists x\, \neg P(x)$
    \item $\forall x\, P(x)$
    \item $\forall x\, \neg P(x)$
    \item $\neg \forall x\, P(x)$
\end{enumerate}

\begin{solnbox}
\textbf{Answer: (C)}

\textbf{Explanation:} De Morgan's Laws for Quantifiers:
$$\neg \exists x\, P(x) \equiv \forall x\, \neg P(x)$$
$$\neg \forall x\, P(x) \equiv \exists x\, \neg P(x)$$

In plain English:
\begin{itemize}[noitemsep]
    \item ``NOT (there exists an $x$ with property $P$)'' = ``ALL $x$ do NOT have property $P$''
    \item If nothing has the property, then everything lacks the property.
\end{itemize}

\textbf{Rule:} Negation flips the quantifier ($\forall \leftrightarrow \exists$) and negates the predicate.
\end{solnbox}

\newpage

%==============================================================================
% PART II: SHORT ANSWER
%==============================================================================
\section*{PART II: Short Answer (20 points)}

\textbf{Question 1.} (5 points) Construct a truth table for $(p \to \neg q) \leftrightarrow (q \lor \neg p)$.

\begin{solnbox}
\textbf{Step-by-step approach:}

First, identify all subexpressions we need:
\begin{enumerate}[noitemsep]
    \item $\neg q$ (negation of $q$)
    \item $\neg p$ (negation of $p$)
    \item $p \to \neg q$ (conditional)
    \item $q \lor \neg p$ (disjunction)
    \item Final: $(p \to \neg q) \leftrightarrow (q \lor \neg p)$
\end{enumerate}

\textbf{Remember:}
\begin{itemize}[noitemsep]
    \item $p \to q$ is FALSE only when $p$ is TRUE and $q$ is FALSE
    \item $p \leftrightarrow q$ is TRUE when both sides have the SAME truth value
\end{itemize}

\begin{center}
\begin{tabular}{cc|cc|cc|c}
\toprule
$p$ & $q$ & $\neg p$ & $\neg q$ & $p \to \neg q$ & $q \lor \neg p$ & $(p \to \neg q) \leftrightarrow (q \lor \neg p)$ \\
\midrule
T & T & F & F & F & T & F \\
T & F & F & T & T & F & F \\
F & T & T & F & T & T & T \\
F & F & T & T & T & T & T \\
\bottomrule
\end{tabular}
\end{center}

\textbf{Detailed calculation for Row 1} ($p=T$, $q=T$):
\begin{itemize}[noitemsep]
    \item $\neg q = \neg T = F$
    \item $\neg p = \neg T = F$
    \item $p \to \neg q = T \to F = F$ (true implies false = false!)
    \item $q \lor \neg p = T \lor F = T$
    \item $F \leftrightarrow T = F$ (different values, so biconditional is false)
\end{itemize}

\textbf{Observation:} This is NOT a tautology (has F values) and NOT a contradiction (has T values). It's a contingency.
\end{solnbox}

\vspace{15pt}

\textbf{Question 2.} (5 points) Translate the following statement into propositional logic using the propositions provided.

\textit{``You can graduate only if you have completed all required courses and you have paid all fees, unless you have a special exemption.''}

Let:
\begin{itemize}[noitemsep]
    \item $g$: ``You can graduate''
    \item $c$: ``You have completed all required courses''
    \item $f$: ``You have paid all fees''
    \item $e$: ``You have a special exemption''
\end{itemize}

\begin{solnbox}
\textbf{Breaking down the English:}

\textbf{Key phrase 1:} ``$A$ only if $B$'' translates to $A \to B$

So ``You can graduate only if [condition]'' means $g \to [\text{condition}]$

\textbf{Key phrase 2:} ``[X] unless [Y]'' translates to $\neg Y \to X$ or equivalently $X \lor Y$

\textbf{Key phrase 3:} ``[A] and [B]'' translates to $A \land B$

\textbf{Putting it together:}

The condition for graduation is: ``completed courses AND paid fees'' = $(c \land f)$

But there's an exception: ``unless you have a special exemption''

This means: Either you meet the requirements $(c \land f)$, OR you have an exemption $e$.

\textbf{Full translation:}
$$\boxed{g \to ((c \land f) \lor e)}$$

\textbf{Alternative interpretation:} If ``unless'' modifies the entire statement:
$$g \to (c \land f) \lor e$$

Both are acceptable depending on how you parse the English. The key insight is recognizing ``only if'' creates an implication with $g$ as the antecedent.
\end{solnbox}

\vspace{15pt}

\textbf{Question 3.} (4 points) Determine the truth value of each statement if the domain for all variables is the set of integers.

\begin{enumerate}[label=(\roman*)]
    \item $\forall n \exists m (n + m = 0)$
    \item $\exists n \forall m (nm = m)$
\end{enumerate}

\begin{solnbox}
\textbf{(i) $\forall n \exists m (n + m = 0)$}

\textbf{Translation:} ``For every integer $n$, there exists an integer $m$ such that $n + m = 0$.''

\textbf{Analysis:} For any integer $n$, we need to find some $m$ where $n + m = 0$.

Simply choose $m = -n$! Then $n + (-n) = 0$. \checkmark

Since every integer has an additive inverse in the integers, this is \textbf{TRUE}.

\rule{\linewidth}{0.3pt}

\textbf{(ii) $\exists n \forall m (nm = m)$}

\textbf{Translation:} ``There exists an integer $n$ such that for ALL integers $m$, $nm = m$.''

\textbf{Analysis:} We need ONE specific $n$ that works for EVERY $m$.

If $nm = m$ for all $m$, then $nm - m = 0$, so $m(n-1) = 0$ for all $m$.

This only works if $n - 1 = 0$, i.e., $n = 1$.

Let's verify: If $n = 1$, then $nm = 1 \cdot m = m$. \checkmark

Yes! The number $n = 1$ (the multiplicative identity) satisfies this.

This is \textbf{TRUE}.
\end{solnbox}

\vspace{15pt}

\textbf{Question 4.} (6 points) Let $F(x, y)$ be the statement ``$x$ is friends with $y$'' where the domain for both $x$ and $y$ is all students at your university. Express each statement using quantifiers.

\begin{enumerate}[label=(\roman*)]
    \item Everyone has at least one friend.
    \item There is a student who is friends with everyone.
    \item No student is friends with everyone.
\end{enumerate}

\begin{solnbox}
\textbf{(i) Everyone has at least one friend.}

``For every student $x$, there exists some student $y$ such that $x$ is friends with $y$.''

$$\boxed{\forall x \exists y\, F(x, y)}$$

\textbf{Note:} The order matters! ``$\forall x \exists y$'' means each person has \textit{their own} friend (possibly different for each person).

\rule{\linewidth}{0.3pt}

\textbf{(ii) There is a student who is friends with everyone.}

``There exists a student $x$ such that for all students $y$, $x$ is friends with $y$.''

$$\boxed{\exists x \forall y\, F(x, y)}$$

\textbf{Key insight:} Here we're looking for ONE specific person who has the universal property. The $\exists$ comes first because we're claiming such a person exists.

\rule{\linewidth}{0.3pt}

\textbf{(iii) No student is friends with everyone.}

This is the NEGATION of (ii)!

``It is NOT the case that there exists a student who is friends with everyone.''

Method 1 (direct): $\neg \exists x \forall y\, F(x, y)$

Method 2 (simplified using De Morgan's):
$$\neg \exists x \forall y\, F(x, y) \equiv \forall x \exists y\, \neg F(x, y)$$

$$\boxed{\forall x \exists y\, \neg F(x, y)}$$

\textbf{In English:} ``For every student, there is someone they are NOT friends with.''
\end{solnbox}

\newpage

%==============================================================================
% PART III: PROBLEM SOLVING
%==============================================================================
\section*{PART III: Problem Solving (25 points)}

\textbf{Question 5.} (6 points) Identify the rule of inference used in each argument below.

\begin{enumerate}[label=(\roman*)]
    \item If it rains, the ground is wet. The ground is not wet. Therefore, it did not rain.
    \item Alice will study math or physics. Alice will not study math. Therefore, Alice will study physics.
    \item If I wake up early, I will exercise. If I exercise, I will be healthy. Therefore, if I wake up early, I will be healthy.
\end{enumerate}

\begin{solnbox}
\textbf{(i) ``If it rains, the ground is wet. The ground is not wet. Therefore, it did not rain.''}

Let $r$ = ``it rains'', $w$ = ``the ground is wet''

Structure:
\begin{align*}
&r \to w \\
&\neg w \\
&\therefore \neg r
\end{align*}

This is \textbf{Modus Tollens}!

\textbf{Why it works:} If the implication $r \to w$ holds, and we observe $\neg w$ (the conclusion is false), then the hypothesis $r$ must have been false. Otherwise, a true $r$ with a true implication would force $w$ to be true --- contradiction!

\rule{\linewidth}{0.3pt}

\textbf{(ii) ``Alice will study math or physics. Alice will not study math. Therefore, Alice will study physics.''}

Let $m$ = ``Alice studies math'', $p$ = ``Alice studies physics''

Structure:
\begin{align*}
&m \lor p \\
&\neg m \\
&\therefore p
\end{align*}

This is \textbf{Disjunctive Syllogism}!

\textbf{Why it works:} In a disjunction, at least one must be true. If we eliminate one option (by knowing it's false), the other must be the true one. Process of elimination!

\rule{\linewidth}{0.3pt}

\textbf{(iii) ``If I wake up early, I will exercise. If I exercise, I will be healthy. Therefore, if I wake up early, I will be healthy.''}

Let $w$ = ``wake up early'', $e$ = ``exercise'', $h$ = ``be healthy''

Structure:
\begin{align*}
&w \to e \\
&e \to h \\
&\therefore w \to h
\end{align*}

This is \textbf{Hypothetical Syllogism}!

\textbf{Why it works:} We chain implications together. If $A$ leads to $B$, and $B$ leads to $C$, then $A$ leads to $C$. We cut out the ``middle man'' $e$.
\end{solnbox}

\vspace{15pt}

\textbf{Question 6.} (7 points) Prove that $\neg(p \to q) \equiv p \land \neg q$ without using truth tables. List each logical law used.

\begin{solnbox}
\textbf{Strategy:} Start with the left side and transform step-by-step until we get the right side.

\textbf{Key insight:} First convert the implication using $p \to q \equiv \neg p \lor q$.

\begin{align*}
\neg(p \to q) &\equiv \neg(\neg p \lor q) && \text{[Implication Law: } p \to q \equiv \neg p \lor q \text{]} \\
&\equiv \neg(\neg p) \land \neg q && \text{[De Morgan's Law: } \neg(A \lor B) \equiv \neg A \land \neg B \text{]} \\
&\equiv p \land \neg q && \text{[Double Negation Law: } \neg(\neg p) \equiv p \text{]}
\end{align*}

$$\boxed{\neg(p \to q) \equiv p \land \neg q}$$

\textbf{What does this mean intuitively?}

The implication $p \to q$ is false ONLY when $p$ is true and $q$ is false.

So the NEGATION of the implication (when is $p \to q$ false?) is exactly when $p$ is true AND $q$ is false, i.e., $p \land \neg q$.

\textbf{Laws used:}
\begin{enumerate}[noitemsep]
    \item Implication Law
    \item De Morgan's Law
    \item Double Negation Law
\end{enumerate}
\end{solnbox}

\vspace{15pt}

\textbf{Question 7.} (4 points) Negate the following statement and simplify. Express your final answer without negation symbols directly in front of quantifiers.

$$\forall x \exists y \forall z\, P(x, y, z)$$

\begin{solnbox}
\textbf{Method:} Push the negation inward, flipping each quantifier as we pass it.

\textbf{Rule:} $\neg \forall$ becomes $\exists \neg$ and $\neg \exists$ becomes $\forall \neg$

\begin{align*}
\neg(\forall x \exists y \forall z\, P(x, y, z)) &\equiv \exists x\, \neg(\exists y \forall z\, P(x, y, z)) && \text{[Flip } \forall \text{ to } \exists \text{]} \\
&\equiv \exists x \forall y\, \neg(\forall z\, P(x, y, z)) && \text{[Flip } \exists \text{ to } \forall \text{]} \\
&\equiv \exists x \forall y \exists z\, \neg P(x, y, z) && \text{[Flip } \forall \text{ to } \exists \text{]}
\end{align*}

$$\boxed{\exists x \forall y \exists z\, \neg P(x, y, z)}$$

\textbf{Pattern recognition:} Notice how each quantifier flips:
\begin{itemize}[noitemsep]
    \item $\forall x \to \exists x$
    \item $\exists y \to \forall y$  
    \item $\forall z \to \exists z$
\end{itemize}

The negation ends up on the predicate $P$ itself.

\textbf{English interpretation:} 
\begin{itemize}[noitemsep]
    \item Original: ``For all $x$, there exists $y$ such that for all $z$, $P$ holds.''
    \item Negation: ``There exists $x$ such that for all $y$, there exists $z$ where $P$ does NOT hold.''
\end{itemize}
\end{solnbox}

\vspace{15pt}

\textbf{Question 8.} (8 points) Prove the following argument is valid by giving a proof using rules of inference. List each rule used.

\textbf{Premises:}
\begin{enumerate}[noitemsep]
    \item $p \lor q$
    \item $p \to r$
    \item $q \to s$
    \item $\neg r$
\end{enumerate}

\textbf{Conclusion:} $s$

\begin{solnbox}
\textbf{Strategy:} We have $p \lor q$ and need to get to $s$. Let's see which path leads to $s$.

From premise 3: $q \to s$. So if we can prove $q$, we get $s$ by Modus Ponens.

How do we get $q$? We have $p \lor q$. If we can show $\neg p$, then by Disjunctive Syllogism, we get $q$!

How do we get $\neg p$? We have $p \to r$ and $\neg r$. By Modus Tollens, we get $\neg p$!

\textbf{Formal Proof:}

\begin{center}
\begin{tabular}{c|l|l}
\toprule
\textbf{Step} & \textbf{Statement} & \textbf{Reason} \\
\midrule
1 & $p \to r$ & Premise \\
2 & $\neg r$ & Premise \\
3 & $\neg p$ & Modus Tollens (1, 2) \\
4 & $p \lor q$ & Premise \\
5 & $q$ & Disjunctive Syllogism (3, 4) \\
6 & $q \to s$ & Premise \\
7 & $s$ & Modus Ponens (5, 6) \\
\bottomrule
\end{tabular}
\end{center}

$$\boxed{\therefore s}$$

\textbf{Proof walkthrough:}
\begin{itemize}[noitemsep]
    \item Steps 1-3: We use $\neg r$ to ``work backwards'' through $p \to r$. Since the conclusion $r$ is false, the hypothesis $p$ must be false. (Modus Tollens)
    \item Steps 4-5: Knowing $\neg p$ and $p \lor q$, we eliminate $p$ from the disjunction, leaving $q$. (Disjunctive Syllogism)
    \item Steps 6-7: Finally, with $q$ and $q \to s$, we conclude $s$. (Modus Ponens)
\end{itemize}
\end{solnbox}

\vspace{20pt}

\begin{center}
\rule{0.5\textwidth}{0.5pt}\\[5pt]
\textbf{END OF EXAM}\\[3pt]
\textit{Total: 55 points}
\end{center}

\end{document}
